\section{The causal completion problem}\label{sec:problem}
A preparation and its legal experimental kernels determine which future tests can actually be performed. They do not supply a manifold, a Fisher matrix, or an observer with an exact posterior register. The question is to connect the predictive information obtained from these tests to a causal implementation whose acquisition controller and memory are both charged. We use \emph{General Theta} for this mathematical programme: prepared kernels, executable tests, a measurable predictive quotient, acquired geometry, certified causal comparison, and decision risk. Generality means a theorem on a stated class, not that every experiment possesses every regularity property.

There are two different obstacles. An arbitrary small terminal code may discard information needed by a later experiment; continuation responses reveal that obstruction. Even when a small code can preserve the relevant response, the observer must decide its next experiment without consulting the discarded history. Strongly separated response refinements resolve the second obstacle by tree pruning, but ordinary noisy reports have overlapping conditional supports. An informative report is not automatically a separated refinement.

Theorem~\ref{thm:pooling} treats the latter case without a support gap. It constructs a causal finite-state observer whose retained labels carry their \emph{actual} posterior distributions. After an action and a report, all incoming posterior values assigned to a cell are pooled together. The cell's read-only representative is their conditional mean under the law induced by this very observer. A Bellman identity then expresses the loss of information exactly as acquired Jensen defects. Concavity controls those defects even where the optimal action switches discontinuously. Lemma~\ref{lem:integrated} proves the required integrated second-order estimate for nonsmooth concave values.

Write $B_n^*$ for the optimal full-history risk with $n$ observations. If every one-step posterior law, uniformly over incoming beliefs and legal actions, has bounded density on a compact interior $d$-dimensional set, then
\begin{equation}\label{eq:intromemory}
 R_\cp(n,M)-B_n^*\asymp M^{-2/d},\qquad
 R_\aut(n,M)-B_n^*\asymp_n M^{-2/d}.
\end{equation}
The first constants are uniform in the horizon when the kernel bounds are. For the second statement the quantitative upper is
\[
 C n\left\lfloor\frac{M-1}{n}\right\rfloor^{-2/d};
\]
the machine has at most $1+n\lfloor(M-1)/n\rfloor\le M$ states, including every phase. Thus \eqref{eq:intromemory} is a fixed-horizon completion law, not a uniform joint $n,M$ equivalence. The full-history acquisition policy is optimized, rather than fixed and subsequently compared with a different control problem.

Theorem~\ref{thm:sensors} verifies the posterior geometry from positive finite-state transitions and continuous report densities by an explicit change of variables. In a two-state specialization, the best last sensor switches at posterior probability $1/2$. A two-call feedback policy has strictly lower risk than every fixed two-sensor sequence. No finite report excludes either hidden state. Theorem~\ref{thm:strata} allows positive-dimensional acquired laws and atoms in one controller, and its validation-sensor realization repeatedly changes between those ranks. These are consequences of the general pooling theorem, not its premises.

The complete separated-acquisition and continuation programmes are supplied as integral Supplement T~\cite{retained}; the earlier block-stability, stopped-task and instrument proofs remain in integral Supplement S~\cite{companion}. No substantive proof has been deleted. Supplement T gives a horizon-uniform completion law on separated response trees, together with distinct noisy continuation obstructions; the present paper closes the overlapping-noise adaptive case on a new class, with the horizon dependence above. We do not claim an intrinsic necessary-and-sufficient classification of all causal experiments, an optimal planning algorithm, or a matched region for every computational resource.

\subsection{Experiments, tests and the state model}
All spaces are standard Borel. A prepared experiment has probability kernels
\[
 K_t^\lambda(dx',dy,dc\mid x,a),
\]
where $c$ records physical calls and time. Policies and simulators are independent of the unknown parameter $\lambda$; all preparation laws, calibration data and policy classes used in a Bayes assertion are specified. A visible history includes actions, reports, publicly supplied costs and legal interface information. An executable future test is a legal finite continuation followed by a bounded scored readout. Histories are predictively equivalent when their legal interfaces and every such expectation agree. Appendix~\ref{sec:quotient} records a measurable realization and its exact minimality condition. A task realization need not simulate the complete experiment.

At a persistent cut an autonomous machine has only its current label in a fixed set of at most $M$ states. Its next action, stop instruction and final readout depend on that label alone; its transition may also use the current report and fresh independent randomness. No elapsed time, program position, earlier output, old report, or private seed remains available unless counted in the state. Read-only instructions may depend on the known experiment and budgets, not the unknown parameter or acquired data. Temporary computation is erased at the next cut. Abstract Borel instructions and known real constants are allowed in the cardinality theorem; Section~\ref{sec:resources} separately gives a certified finite-precision implementation. Existence of the programme is not efficient synthesis of it.

A fixed $n$-call task permits actions at each of the $n$ stages and reveals the scored audit only after the forecast. This theorem does not optimize the physical stopping time: its phase is internal to the finite machine. Bounded stopping with a nonintervening audit has a separate result in Supplement T; no fixed-time statement is substituted at an arbitrary stopping time here. A checkpoint procedure may use the full history during acquisition and retain only $M$ labels at its terminal cut. Both procedure classes use the same preparation, actions and audit.

For a finite measure $\nu$ in a Euclidean task space set
\begin{equation}\label{eq:quantization}
 q_M(\nu)=\inf_{c_1,\ldots,c_M}\int\min_j\norm{x-c_j}^2\,\nu(dx).
\end{equation}
Subprobability mass is not normalized away. For a fixed deterministic acquisition policy $\pi$, let $P_n^\pi=\E[Y\mid H_n^\pi]$, $B_n^\pi=\E\norm{Y-P_n^\pi}^2$, and $\nu_n^\pi=\law(P_n^\pi)$. Projection gives the fixed-policy identity
\begin{equation}\label{eq:fixedcp}
 \inf_{M\text{-label terminal codes}}\E\norm{Y-\widehat Y}^2
       =B_n^\pi+q_M(\nu_n^\pi).
\end{equation}
The encoder assigns the least-index nearest center; such ties are Borel. An infimum suffices if an optimal codebook is not attained.

\begin{lemma}[Independent randomization at a checkpoint]\label{lem:seeds}
Suppose all acquisition and coding coins are independent of the prepared physical randomness. Allowing a full-history policy and its checkpoint code to use these coins does not lower the infimum in \eqref{eq:fixedcp} after joint optimization over deterministic policies and codes. In particular,
\begin{equation}\label{eq:checkpointfunctional}
 R_\cp(n,M)=\inf_{\pi\ {\rm deterministic}}
             \{B_n^\pi+q_M(\nu_n^\pi)\}.
\end{equation}
\end{lemma}
\begin{proof}
Supply the encoder and decoder with the entire independent seed, as a relaxation if necessary. For each fixed seed their full-history policy and code are deterministic; their risk is bounded below by the deterministic joint infimum. Integrate over the seed. Deterministic procedures are admitted in the randomized class, proving equality. This argument concerns the full-history checkpoint problem. Fixing a time-indexed coin tape need not preserve the autonomous state budget, and no such purification is asserted. For an autonomous squared-loss lower bound one may first replace each randomized terminal readout by its conditional mean, obtaining at most $M$ fixed centers, and then relax its earlier control to full-history control.
\end{proof}
The policy infimum in \eqref{eq:checkpointfunctional} is not interchanged with either of its policy-dependent summands.
